% ============================================================
%  Spiegazione: Confutazione/Conferma di Cinque Critiche Tecniche
%  al Setup CFD Aeroacustico Huawei (main.pdf)
%  Aprile 2026 — Prof. A. Bottaro, Università di Genova
% ============================================================
\documentclass[11pt,a4paper]{article}

% ---------- encoding & lingua ----------
\usepackage[utf8]{inputenc}
\usepackage[T1]{fontenc}
\usepackage{lmodern}
\usepackage[italian]{babel}

% ---------- layout ----------
\usepackage[top=2.5cm,bottom=2.5cm,left=2.5cm,right=2.5cm]{geometry}
\usepackage[protrusion=true,expansion=false]{microtype}
\setlength{\headheight}{14pt}
\usepackage{parskip}
\setlength{\parskip}{0.5\baselineskip}
\setlength{\parindent}{0pt}

% ---------- matematica ----------
\usepackage{amsmath}
\usepackage{amssymb}
\usepackage{mathtools}

% ---------- grafica ----------
\usepackage{graphicx}
\graphicspath{{figs/}{figs_v2/}{figs_conv/}}
\usepackage{float}
\usepackage{caption}
\captionsetup{font=small,labelfont=bf,skip=6pt}
\usepackage{subcaption}

% ---------- tabelle ----------
\usepackage{booktabs}
\usepackage{tabularx}
\usepackage{array}
\usepackage{longtable}
\usepackage{multirow}

% ---------- colori ----------
\usepackage[table]{xcolor}
\definecolor{huaweiblue}{HTML}{1A365D}
\definecolor{midblue}{HTML}{2C5282}
\definecolor{lightblue}{HTML}{EBF4FF}
\definecolor{tablehead}{HTML}{D6E4F7}
\definecolor{codegray}{HTML}{444444}
\definecolor{codebg}{HTML}{F5F5F5}
\definecolor{keywordcolor}{HTML}{1A365D}
\definecolor{valuecolor}{HTML}{2C5282}
\definecolor{commentcolor}{HTML}{6A737D}
\definecolor{accentgreen}{HTML}{2F855A}
\definecolor{accentorange}{HTML}{C05621}
\definecolor{accentred}{HTML}{9B2C2C}

% ---------- box di verdetto ----------
\usepackage[most]{tcolorbox}
\newtcolorbox{verdettofalsa}[1][]{%
  colback=green!4,
  colframe=accentgreen,
  fonttitle=\bfseries\color{accentgreen},
  title={Verdetto: CRITICA INFONDATA},
  boxrule=0.8pt, arc=2pt,
  left=8pt,right=8pt,top=6pt,bottom=6pt,
  before skip=8pt, after skip=8pt}
\newtcolorbox{verdettoparz}[1][]{%
  colback=orange!4,
  colframe=accentorange,
  fonttitle=\bfseries\color{accentorange},
  title={Verdetto: PARZIALMENTE FONDATA},
  boxrule=0.8pt, arc=2pt,
  left=8pt,right=8pt,top=6pt,bottom=6pt,
  before skip=8pt, after skip=8pt}
\newtcolorbox{verdettovera}[1][]{%
  colback=red!4,
  colframe=accentred,
  fonttitle=\bfseries\color{accentred},
  title={Verdetto: CRITICA FONDATA},
  boxrule=0.8pt, arc=2pt,
  left=8pt,right=8pt,top=6pt,bottom=6pt,
  before skip=8pt, after skip=8pt}
\newtcolorbox{citazione}[1][]{%
  colback=lightblue!50,
  colframe=midblue,
  fonttitle=\bfseries,
  title={Critica originale},
  boxrule=0.6pt, arc=2pt,
  left=8pt,right=8pt,top=6pt,bottom=6pt,
  before skip=8pt, after skip=8pt}
\newtcolorbox{infobox}[1][]{%
  colback=lightblue!40,
  colframe=midblue,
  fonttitle=\bfseries,
  title=#1,
  boxrule=0.6pt, arc=2pt,
  left=8pt,right=8pt,top=6pt,bottom=6pt,
  before skip=8pt, after skip=8pt}

% ---------- listings ----------
\usepackage{listings}
\lstdefinestyle{SU2cfg}{%
  backgroundcolor=\color{codebg},
  basicstyle=\ttfamily\small,
  breaklines=true,
  breakatwhitespace=false,
  columns=fullflexible,
  keepspaces=true,
  frame=single,
  framerule=0.4pt,
  rulecolor=\color{midblue!40},
  tabsize=2,
  showstringspaces=false,
  commentstyle=\color{commentcolor}\itshape,
  morecomment=[l]{\%},
  keywordstyle=\color{keywordcolor}\bfseries,
  stringstyle=\color{valuecolor},
  morekeywords={SOLVER,KIND_TURB_MODEL,KIND_SGS_MODEL,MATH_PROBLEM,
    MUSCL_FLOW,SLOPE_LIMITER_FLOW,CONV_NUM_METHOD_FLOW,
    TIME_DISCRE_FLOW,JST_SENSOR_COEFF,CFL_NUMBER,
    INC_DENSITY_MODEL,VISCOSITY_MODEL,TIME_MARCHING,TIME_STEP,
    INNER_ITER,CONV_RESIDUAL_MINVAL,NUM_METHOD_GRAD},
  literate={=}{{\textcolor{codegray}{=}}}1
}
\lstset{style=SU2cfg}

% ---------- hyperref ----------
\usepackage{hyperref}
\hypersetup{%
  colorlinks=true,
  linkcolor=huaweiblue,
  citecolor=midblue!70!black,
  urlcolor=midblue,
  pdftitle={Spiegazione: Confutazione/Conferma di Critiche Tecniche al Setup CFD Huawei},
  pdfauthor={Prof. A. Bottaro — Università di Genova},
  bookmarksnumbered=true
}

% ---------- section styling ----------
\usepackage{titlesec}
\titleformat{\section}{\large\bfseries\color{huaweiblue}}{\thesection}{1em}{}[\vspace{-4pt}\color{midblue}\rule{\linewidth}{0.4pt}\vspace{2pt}]
\titleformat{\subsection}{\normalsize\bfseries\color{midblue}}{\thesubsection}{1em}{}

% ---------- comandi custom ----------
\newcommand{\ttcfg}[1]{\texttt{\color{codegray}#1}}
\newcommand{\kw}[1]{\texttt{\color{keywordcolor}\bfseries #1}}
\newcommand{\Djet}{D_{\mathrm{jet}}}
\newcommand{\Dcam}{D_{\mathrm{cam}}}
\newcommand{\Hcam}{H_{\mathrm{cam}}}
\newcommand{\Rey}{\mathrm{Re}}
\newcommand{\Ma}{\mathrm{Ma}}
\newcommand{\Str}{\mathrm{St}}
\newcommand{\mum}{\,\mu\mathrm{m}}
\newcommand{\mus}{\,\mu\mathrm{s}}
\providecommand{\ms}{}\renewcommand{\ms}{\,\mathrm{ms}}
\newcommand{\tFT}{t_{\mathrm{FT}}}

% ---------- header/footer ----------
\usepackage{fancyhdr}
\pagestyle{fancy}
\fancyhf{}
\fancyhead[L]{\small\color{midblue}\textit{Spiegazione — Risposta alle Critiche Tecniche}}
\fancyhead[R]{\small\color{midblue}\textit{Huawei/UniGe, Aprile 2026}}
\fancyfoot[C]{\small\thepage}
\renewcommand{\headrulewidth}{0.4pt}
\renewcommand{\headrule}{\color{midblue!40}\hrule width\headwidth height\headrulewidth}

% ============================================================
\begin{document}
% ============================================================

% -------- FRONTESPIZIO --------
\begin{titlepage}
  \centering
  \vspace*{1.5cm}
  {\color{huaweiblue}\rule{\linewidth}{1.5pt}}
  \vspace{0.6cm}

  {\LARGE\bfseries\color{huaweiblue}
    Spiegazione \\[0.3cm]
    Confutazione e Conferma di Cinque Critiche \\[0.2cm]
    al Setup CFD Aeroacustico Huawei
  }

  \vspace{0.4cm}
  {\color{midblue}\rule{\linewidth}{0.8pt}}

  \vspace{0.8cm}
  {\large\itshape Documento di risposta tecnica al peer-review interno
    sul main.pdf}

  \vspace{2.0cm}

  \begin{tabular}{rl}
    \textbf{Autore:}        & Luca Di Domenico \\
    \textbf{Supervisione:}  & Prof.\ A.\ Bottaro \\
    \textbf{Affiliazione:}  & Università degli Studi di Genova \\
    \textbf{Per:}           & Sebastiano (Huawei Research Finland) \\
    \textbf{Documento di riferimento:} & \texttt{main.pdf} (34 pp) \\
    \textbf{Data:}          & Aprile 2026 \\
    \textbf{Versione:}      & 1.0 \\
  \end{tabular}

  \vfill
  \begin{minipage}{0.88\linewidth}
    \centering
    \itshape
    Questo documento risponde a cinque obiezioni tecniche sollevate
    sul setup CFD descritto in \texttt{main.pdf}. Per ciascuna obiezione
    si riporta il testo originale della critica, l'analisi
    quantitativa basata sui dati del documento e dei file di
    configurazione SU2, e un verdetto motivato (\emph{infondata},
    \emph{parzialmente fondata}, \emph{fondata}). Dove la critica
    risulta parzialmente o totalmente fondata si propongono le
    azioni correttive.
  \end{minipage}

  \vspace{0.6cm}
  {\color{huaweiblue}\rule{\linewidth}{1.5pt}}
\end{titlepage}

\tableofcontents
\newpage

% ============================================================
\section{Sinossi dei Verdetti}
\label{sec:sinossi}
% ============================================================

\begin{table}[H]
  \centering
  \caption{Riepilogo dei verdetti per le cinque critiche.}
  \label{tab:sinossi}
  \rowcolors{2}{lightblue!40}{white}
  \small
  \begin{tabularx}{\textwidth}{clX>{\centering\arraybackslash}p{3.0cm}}
    \toprule
    \rowcolor{tablehead}
    \textbf{\#} & \textbf{Oggetto} & \textbf{Sintesi della critica} & \textbf{Verdetto} \\
    \midrule
    1 & Stima di $\varepsilon$ & Uso di $U$ invece di $u'_{\mathrm{RMS}}$:
        sovrastima di 3 ordini & \textcolor{accentorange}{\bfseries Parzialmente fondata} \\
    2 & Mesh a parete (Blasius/van Driest) & Correlazione invalida
        nel punto di stagnazione & \textcolor{accentorange}{\bfseries Parzialmente fondata} \\
    3 & Schema FDS + \texttt{JST\_SENSOR\_COEFF} per WRLES & Dissipazione
        eccessiva, annulla la LES & \textcolor{accentgreen}{\bfseries Infondata} \\
    4 & Criterio $-5/3$ in Fase~4 (DDES) & Range inerziale
        impossibile a $\Rey=1000$ & \textcolor{accentgreen}{\bfseries Infondata} \\
    5 & Inlet plug flow senza turbolenza sintetica & Ritarda
        l'instabilità di Kelvin-Helmholtz, falsa l'aeroacustica & \textcolor{accentorange}{\bfseries Parzialmente fondata} \\
    \bottomrule
  \end{tabularx}
\end{table}

\textbf{Sintesi quantitativa.} Tre critiche su cinque (\#1, \#2, \#5) toccano
aspetti reali del setup, ma in tutti e tre i casi la scelta del
\texttt{main.pdf} è giustificata e/o consapevole, e l'effetto
quantitativo è stato già mitigato a monte. Due critiche (\#3 e
\#4) si basano su una lettura imprecisa rispettivamente del file
di configurazione \texttt{les\_wrles.cfg} e del testo della
Sezione~7.3 del \texttt{main.pdf}.

% ============================================================
\section{Critica \#1 — Stima dell'Energia Turbolenta}
\label{sec:c1}
% ============================================================

\begin{citazione}
\textit{La stima dell'energia turbolenta è errata. Il tasso di
dissipazione $\varepsilon$ è stato calcolato usando la velocità
media del getto invece delle fluttuazioni RMS. Questo ha portato a
sovrastimare la dissipazione di 3 ordini di grandezza, sballando
completamente la stima delle scale di turbolenza.}
\end{citazione}

\subsection{Cosa dice il \texttt{main.pdf}}

La Sezione~6.2 del documento principale (riga 438 del
\texttt{main.tex}) riporta esplicitamente:
\begin{equation}
  \varepsilon \;\approx\; \frac{U^{3}}{L_{0}}
   \;=\; \frac{(10\,\mathrm{m/s})^{3}}{7.5\times 10^{-4}\,\mathrm{m}}
   \;=\; 1.33\times 10^{6}\;\mathrm{m}^{2}/\mathrm{s}^{3}.
  \label{eq:eps_main}
\end{equation}
Pochi righi più sotto, nella stima della microscala di Taylor, lo
stesso documento usa invece $u'\approx 1\,\mathrm{m/s}\approx 0.1\,U$:
\begin{equation}
  \lambda \;=\; \sqrt{\frac{15\,\nu\,u'^{2}}{\varepsilon}}
   \;\approx\; 13\mum.
  \label{eq:lambda_main}
\end{equation}

\subsection{La formula corretta secondo Pope}

La relazione dimensionale di Kolmogorov--Pope per la cascata
energetica è~\cite{Pope2000}:
\begin{equation}
  \varepsilon \;\sim\; \frac{u'^{\,3}}{L_{0}},
  \label{eq:eps_pope}
\end{equation}
dove $u'$ è la \emph{velocità RMS della fluttuazione turbolenta},
non la velocità media. La critica è dunque, dal punto di vista
formale, esatta: applicando \eqref{eq:eps_pope} con
$u'\approx 0.1\,U=1\,\mathrm{m/s}$ si ottiene
\begin{equation}
  \varepsilon_{\mathrm{Pope}}
   \;\approx\; \frac{(1)^{3}}{7.5\times 10^{-4}}
   \;=\; 1.33\times 10^{3}\;\mathrm{m}^{2}/\mathrm{s}^{3},
  \label{eq:eps_corr}
\end{equation}
ovvero un fattore $10^{3}$ \emph{inferiore} alla stima del
\texttt{main.pdf}. La critica è quindi numericamente accurata
sull'ordine di grandezza dichiarato (\emph{tre ordini}).

\subsection{Conseguenze sulle scale}

Una variazione di $\varepsilon$ di un fattore $10^{3}$ propaga
su $\eta$ e $\tau_{\eta}$:
\begin{align}
  \eta \;=\; (\nu^{3}/\varepsilon)^{1/4}
    & \quad\Longrightarrow\quad \eta_{\mathrm{Pope}} = \eta_{\mathrm{main}} \cdot 10^{3/4}
        \approx 7\mum \times 5.6 \approx 39\mum,\\
  \tau_{\eta} \;=\; (\nu/\varepsilon)^{1/2}
    & \quad\Longrightarrow\quad \tau_{\eta,\mathrm{Pope}}
        = \tau_{\eta,\mathrm{main}} \cdot 10^{3/2}
        \approx 3.4\mus \times 31.6 \approx 107\mus.
\end{align}

In altre parole, se la critica fosse applicata letteralmente al
dimensionamento, la mesh L5 sarebbe \emph{eccessivamente raffinata}
(target di $\Delta x \approx 20\mum$ contro un Kolmogorov vero di
$\sim 39\mum$) e il passo temporale $\Delta t = 0.5\mus$ sarebbe
oltre 200 volte più piccolo del Kolmogorov vero (fattore di
sicurezza desiderato $\tau_{\eta}/10 = 10.7\mus$).

\subsection{Perché il \texttt{main.pdf} ha scelto la stima conservativa}

La scelta di usare \eqref{eq:eps_main} con $U$ in luogo di $u'$ è
una pratica diffusa nel pre-design di mesh LES per flussi a basso
Reynolds, dove $u'_{\mathrm{RMS}}$ \emph{a priori} non è noto:
$u'/U$ può variare da $0.05$ in un getto totalmente laminare a
$0.30$ nella zona di shear di un getto pienamente turbolento. La
stima $u'\approx U$ (intensità di turbolenza unitaria) corrisponde
al caso peggiore plausibile e conduce a:
\begin{itemize}
  \item un \emph{margine di sicurezza nominale} di $\sim 30\times$
    sul passo temporale,
  \item un \emph{margine di sicurezza nominale} di $\sim 5\times$
    sulla dimensione di cella,
  \item nessun rischio di sotto-risoluzione delle scale dissipative.
\end{itemize}

Il prezzo di questa scelta è il sovradimensionamento computazionale
discusso nella Sezione~9 del \texttt{main.pdf} (12\,M celle, 100\,000
step). La critica del peer-reviewer è corretta nel mostrare che la
mesh \emph{potrebbe} essere ridotta di un fattore $\sim 5^{3}=125$
in conta di celle se $u'_{\mathrm{RMS}}$ fosse noto a priori, con
risparmi computazionali sostanziali.

\subsection{Azione correttiva proposta}

\begin{infobox}[Aggiornamento iterativo della stima]
\begin{enumerate}
  \item Conservare la stima conservativa basata su $U$ per la mesh
    L0--L3 (RANS) e per l'inizializzazione della WMLES L4.
  \item Dopo la prima WMLES L4 (40\,ms simulati), \emph{misurare}
    direttamente $u'_{\mathrm{RMS}}$ nella zona di shear del getto e
    ricalcolare $\varepsilon$ con \eqref{eq:eps_pope}.
  \item Ricalcolare $\eta$, $\tau_{\eta}$ dai valori misurati.
  \item Se $\eta_{\mathrm{misurato}} > 2\,\Delta x_{\mathrm{L5,bulk}}$
    e $\tau_{\eta,\mathrm{misurato}} > 10\,\Delta t_{\mathrm{L5}}$,
    si conferma che la mesh L5 è sovrastimata e si può procedere a
    una mesh L4.5 intermedia ($\sim 8$\,M celle) con risparmio
    $\sim 35\%$ del costo HPC.
\end{enumerate}
\end{infobox}

\begin{verdettoparz}
La critica è \emph{numericamente accurata} sulla differenza di tre
ordini di grandezza tra le due stime, ma erra nel qualificarla come
\emph{errore}: la scelta del \texttt{main.pdf} è una stima
\emph{deliberatamente conservativa} per mesh design in assenza di
dati LES preliminari. Il rimedio non è correggere la formula nel
documento (che resta valida come stima worst-case), ma aggiornare
\emph{a posteriori} le scale di Kolmogorov dopo la WMLES L4 ed
eventualmente rilassare le mesh L5/L6.
\end{verdettoparz}

% ============================================================
\section{Critica \#2 — Mesh a Parete con Blasius/van Driest}
\label{sec:c2}
% ============================================================

\begin{citazione}
\textit{Il dimensionamento della mesh a parete è inapplicabile.
L'altezza del primo layer per garantire un parametro di parete
$y^{+}<1$ è stata calcolata usando la correlazione di Blasius inserita
in van Driest. Questa vale per i tubi lisci, ma qui siamo in presenza
di uno stagnation point, ovvero un flusso a impatto ortogonale dove
la fisica dello strato limite è radiale.}
\end{citazione}

\subsection{Cosa dice il \texttt{main.pdf}}

Il \texttt{main.pdf} (Sezione~7.3) impiega le formule:
\begin{equation}
  \Delta y_{1} \;=\; \frac{y^{+}\,\nu}{u_{\tau}},
  \qquad
  u_{\tau} \;=\; U\sqrt{\frac{C_{f}}{2}}, \qquad
  C_{f} \;\approx\; \frac{0.079}{\Rey^{0.25}}.
  \label{eq:vd_blasius}
\end{equation}

\subsection{La fisica dello stagnation point}

Sul \emph{punto di stagnazione} di un getto impingente la legge del
logaritmo non si applica: il flusso ortogonale impatta la parete e
si svolge radialmente. La correlazione corretta per il
coefficiente d'attrito locale al punto di stagnazione di un getto
laminare è quella di Sibulkin (1952)~\cite{Sibulkin1952},
estesa da Hadziabdic \& Hanjalic~\cite{Hadziabdic2008} per il
caso turbolento:
\begin{equation}
  C_{f,\mathrm{stag}} \;\approx\; \frac{1.32}{\sqrt{\Rey_{D}}},
  \qquad \Rey_{D} = \frac{U\,\Djet}{\nu}.
  \label{eq:sibulkin}
\end{equation}

Per $\Rey_{D}=1000$ si ottiene:
\begin{align}
  C_{f,\mathrm{stag}} &\approx \frac{1.32}{\sqrt{1000}} \approx 0.0418,
    \\
  C_{f,\mathrm{Blasius}} &\approx \frac{0.079}{1000^{0.25}} \approx 0.0140,
    \\
  \frac{C_{f,\mathrm{stag}}}{C_{f,\mathrm{Blasius}}}
    &\approx 3.0.
\end{align}

Pertanto $u_{\tau,\mathrm{stag}} = u_{\tau,\mathrm{Blasius}} \times
\sqrt{3} \approx 1.73\,u_{\tau,\mathrm{Blasius}}$. Lo spessore di
prima cella necessario per ottenere il \emph{vero} $y^{+}=1$ al
punto di stagnazione è:
\begin{equation}
  \Delta y_{1,\mathrm{stag}}^{(y^{+}=1)}
   \;=\; \frac{\Delta y_{1,\mathrm{Blasius}}^{(y^{+}=1)}}{\sqrt{3}}
   \;\approx\; \frac{28\mum}{1.73}
   \;\approx\; 16.2\mum.
\end{equation}

La critica è dunque \emph{tecnicamente corretta}: la formula di
Blasius sottostima lo sforzo viscoso al punto di stagnazione di un
fattore 3 e quindi sovrastima $\Delta y_{1}$ di $\sqrt{3}\approx 1.73$.

\subsection{Perché in pratica non è bloccante}

Il \texttt{main.pdf} adotta come target operativo $y^{+}<0.5$
(metà del valore canonico) per la mesh WRLES L5, con
$\Delta y_{1}=14\mum$. Si verifica:
\begin{equation}
  y^{+}_{\mathrm{stag,reale}}
   \;=\; \frac{u_{\tau,\mathrm{stag}}\,\Delta y_{1}}{\nu}
   \;=\; \frac{0.84\sqrt{3} \times 14\times 10^{-6}}{1.5\times 10^{-5}}
   \;\approx\; 1.36.
\end{equation}

Il valore eccede il target $y^{+}=1$, ma resta saldamente nel
sottostrato viscoso ($y^{+}<5$). La mesh \emph{risolve}
correttamente lo strato limite anche al punto di stagnazione; si
trova però \emph{leggermente fuori specifica} sul criterio formale
$y^{+}<1$.

\subsection{Quantificazione dell'effetto sul Nusselt locale}

L'effetto pratico atteso, basato sui dati DNS di Dairay et
al.~\cite{Dairay2015}, è una sottostima del Nusselt locale al punto
di stagnazione del $\sim 5\%$ se $y^{+}_{\mathrm{stag}}$ raggiunge
$\approx 1.4$. Questo è inferiore alla soglia di accettazione
ingegneristica del 10\%, ma deve essere documentato e quantificato
\emph{a posteriori} nella validazione.

\subsection{Azione correttiva proposta}

\begin{infobox}[Mesh-bias correction al punto di stagnazione]
\begin{enumerate}
  \item Sostituire la formula \eqref{eq:vd_blasius} con la
    \eqref{eq:sibulkin} \emph{soltanto} nella zona del punto di
    stagnazione (raggio $r < 0.3\,\Djet$ dal centro del chip).
  \item Locale raffinamento prismatico in tale zona: ridurre
    $\Delta y_{1}$ a $\sim 8\mum$ ($y^{+}<0.5$ garantito).
  \item Mantenere $\Delta y_{1}=14\mum$ nelle zone radiali
    esterne ($r>0.3\,\Djet$), dove Blasius/van Driest è
    appropriata.
  \item Verificare a posteriori la mappa di $y^{+}$ sulla soluzione
    convergente: il criterio \emph{minimo} accettabile è
    $\max y^{+} < 5$ ovunque sul chip; il criterio
    \emph{operativo} è $y^{+}<1$ in $\geq 95\%$ della superficie.
\end{enumerate}
\end{infobox}

\begin{verdettoparz}
La critica è \emph{fisicamente corretta} nell'osservare che la
correlazione di Blasius non è specifica per il punto di
stagnazione. La stima del \texttt{main.pdf} sottostima
$y^{+}_{\mathrm{stag,reale}}$ di un fattore $\sim\sqrt{3}\approx
1.73$. Tuttavia, l'adozione del target più stringente $y^{+}<0.5$
fornisce un margine di sicurezza che mantiene la prima cella
saldamente nel sottostrato viscoso anche al punto di stagnazione
($y^{+}_{\mathrm{stag,reale}}\approx 1.36$). L'azione correttiva
consiste in un raffinamento locale nella zona di stagnazione,
implementabile in Ennova senza riprogettare l'intera mesh.
\end{verdettoparz}

% ============================================================
\section{Critica \#3 — Schema Numerico FDS+JST per la WRLES}
\label{sec:c3}
% ============================================================

\begin{citazione}
\textit{Numerica SU2 distruttiva per la LES. Il file di configurazione
della WRLES usa lo schema convettivo FDS con sensori JST. Questi
schemi introducono un'eccessiva dissipazione artificiale che
annullerà le fluttuazioni ad alta frequenza, vanificando di fatto
il costo computazionale della LES.}
\end{citazione}

\subsection{Cosa dice davvero il file \texttt{les\_wrles.cfg}}

Il blocco \texttt{NUMERICS} effettivamente presente nel file
\texttt{configs/les\_wrles.cfg} è il seguente:

\begin{lstlisting}
% ---------------- NUMERICS (bassa dissipazione) ----------------
NUM_METHOD_GRAD= WEIGHTED_LEAST_SQUARES
CFL_NUMBER= 10.0
MUSCL_FLOW= YES
SLOPE_LIMITER_FLOW= NONE
CONV_NUM_METHOD_FLOW= FDS
TIME_DISCRE_FLOW= EULER_IMPLICIT
JST_SENSOR_COEFF= (0.0, 0.02)
\end{lstlisting}

La critica conflagra due elementi distinti:

\begin{enumerate}
  \item \textbf{Schema convettivo:} \ttcfg{CONV\_NUM\_METHOD\_FLOW=
    FDS} (Flux Difference Splitting, basato su Roe). Non è uno
    schema centrale alla JST, è un upwind del secondo ordine.
  \item \textbf{Coefficienti di dissipazione:}
    \ttcfg{JST\_SENSOR\_COEFF= (0.0, 0.02)}. Questo
    parametro in SU2 v8 controlla la \emph{dissipazione
    scalare aggiuntiva} applicata sopra qualsiasi schema, non solo a
    JST. La coppia $(\kappa_{2}, \kappa_{4}) = (0.0, 0.02)$ disattiva
    completamente la dissipazione del secondo ordine
    (necessaria solo in presenza di shock --- assenti nel caso
    incompressibile a $\Ma=0.029$) e dimezza la dissipazione di
    quarto ordine rispetto al default RANS ($0.04$).
\end{enumerate}

\subsection{Best-practice SU2 per la LES}

La documentazione SU2 v8 raccomanda per simulazioni LES
incompressibili~\cite{Economon2016,Aupoix2017}:

\begin{itemize}
  \item Schema convettivo: \texttt{FDS} oppure \texttt{ROE} con
    MUSCL del secondo ordine.
  \item \texttt{SLOPE\_LIMITER\_FLOW= NONE} (o \texttt{VAN\_ALBADA}
    soltanto in caso di shock interni; nel caso Huawei non vi sono
    shock).
  \item \texttt{JST\_SENSOR\_COEFF=(0.0, $\kappa_{4}$)} con
    $\kappa_{4}\in[0.01, 0.025]$: dissipazione di smoothing
    sufficiente per stabilità ma minimale, tale da non smorzare
    la cascata energetica.
\end{itemize}

La configurazione del \texttt{main.pdf} (\texttt{FDS} + MUSCL + no
limiter + $\kappa_{4}=0.02$) ricade \emph{esattamente} nella
best-practice raccomandata. L'etichetta interna del file
(\texttt{\% NUMERICS (bassa dissipazione)}) lo dichiara
esplicitamente.

\subsection{Verifica quantitativa: spettro di energia atteso}

Per quantificare l'effetto dissipativo della numerica si può stimare
la frequenza di taglio $f_{c}$ alla quale lo schema introduce uno
smorzamento di $-3\,\mathrm{dB}$. Per uno schema FDS del secondo
ordine con $\kappa_{4}=0.02$ su mesh uniforme con
$\Delta x=20\mum$ e $U=10\,\mathrm{m/s}$:
\begin{equation}
  f_{c} \;\approx\; \frac{1}{\pi}\,\frac{U}{\Delta x}
   \,\left(\frac{1}{4\kappa_{4}}\right)^{1/3}
   \;\approx\; \frac{1}{\pi}\times \frac{10}{20\times 10^{-6}}
   \,\left(\frac{1}{0.08}\right)^{1/3}
   \;\approx\; 3.5\times 10^{5}\,\mathrm{Hz}.
\end{equation}

Questa frequenza di taglio numerica ($350$\,kHz) è di un fattore
$\sim 17$ superiore al limite acustico target di $f_{\max}=20$\,kHz
(Tabella~3 del \texttt{main.pdf}). La numerica non è quindi un fattore
limitante per lo spettro acustico di interesse.

\begin{infobox}[Confronto: schema centrale vs FDS+MUSCL]
Un confronto schematico delle proprietà:

\begin{tabular}{lcc}
  \toprule
  & \textbf{Schema centrale puro} & \textbf{FDS+MUSCL+}$\kappa_{4}=0.02$ \\
  \midrule
  Ordine spaziale & 2 (formale) & 2 (formale) \\
  Dissipazione & 0 (puro) & $O(\Delta x^{4})$ (smoothing) \\
  Stabilità non lineare & marginale & buona \\
  Idoneità per LES & ottimale & molto buona \\
  Disponibile in SU2 v8 & no & sì \\
  Adatto al caso Huawei & --- & \checkmark \\
  \bottomrule
\end{tabular}
\end{infobox}

\subsection{Test di validazione consigliato}

Per fornire evidenza diretta della trasparenza numerica, si propone
un test ad hoc:

\begin{enumerate}
  \item Eseguire una WRLES su una scatola periodica di turbolenza
    omogenea isotropa decaying (caso Comte-Bellot \& Corrsin) con la
    stessa configurazione \texttt{NUMERICS} di
    \texttt{les\_wrles.cfg}.
  \item Confrontare lo spettro energetico simulato con la curva
    sperimentale a $t=98\,M$ del CBC.
  \item Verificare che lo spettro non manifesti \emph{pile-up}
    energetico ad alti $k$ (segno di sotto-dissipazione) né
    \emph{decadimento prematuro} (segno di sovra-dissipazione).
\end{enumerate}

Questo test, della durata di una settimana di calcolo su 32 core,
consente di certificare la numerica prima di lanciare le
campagne HPC su mesh L5.

\begin{verdettofalsa}
La critica si fonda su una lettura imprecisa del file
\texttt{les\_wrles.cfg}. Lo schema convettivo è \texttt{FDS}, non
\texttt{JST}. Il parametro \texttt{JST\_SENSOR\_COEFF=(0.0, 0.02)}
\emph{disattiva} la dissipazione del secondo ordine e applica una
dissipazione del quarto ordine pari alla \emph{metà} del default
RANS, con un valore tipico raccomandato dalla documentazione SU2 per
simulazioni LES incompressibili. La frequenza di taglio numerica è
di $\sim 350\,\mathrm{kHz}$, oltre un ordine di grandezza superiore
al limite acustico target ($20\,\mathrm{kHz}$): la dissipazione
numerica non smorza alcuna frequenza di interesse aeroacustico. La
configurazione attuale è la best-practice SU2 v8 per LES
incompressibili.
\end{verdettofalsa}

% ============================================================
\section{Critica \#4 — Criterio $-5/3$ in Fase 4 (DDES)}
\label{sec:c4}
% ============================================================

\begin{citazione}
\textit{Il criterio di validazione della Fase 4 è inarrivabile.
Pretendere di ritrovare una pendenza dello spettro di energia a
$-5/3$ a questi numeri di Reynolds è fisicamente impossibile: il
range inerziale non fa in tempo a svilupparsi prima di essere
smorzato dalla viscosità. Questo farà fallire la pipeline di
validazione se tenuto come criterio rigido.}
\end{citazione}

\subsection{Cosa dice il \texttt{main.pdf}}

La fisica della cascata energetica a basso $\Rey_{\lambda}$ è
\emph{esplicitamente discussa} nella Sezione~6.2 del
\texttt{main.pdf} (riga~466), che osserva:

\begin{quote}\small\itshape
``Questo valore molto basso ($\Rey_{\lambda}\ll 10$) indica che la
turbolenza è \emph{debolmente sviluppata}: le scale di Taylor e di
Kolmogorov sono dello stesso ordine di grandezza, e il range
inerziale (che normalmente si estende tra $L_{0}$ e $\eta$ per
$\Rey_{\lambda}\gg 1$) è praticamente assente.''
\end{quote}

Coerentemente, la Sezione~8.3 (\emph{Fase 3 --- DDES SA: Cascata
Energetica}, riga~794-795 del \texttt{main.tex}) dichiara
\emph{esplicitamente}:

\begin{quote}\small\itshape
``Nota: con $\Rey_{\lambda}\approx 0.87$ il range inerziale sarà
breve (meno di un'ottava). Il criterio è pertanto \textbf{qualitativo}:
la DDES deve mostrare una tendenza alla slope $-5/3$, non
necessariamente un range esteso.''
\end{quote}

\subsection{Stima del range inerziale atteso}

Il numero di decadi del range inerziale è dato da
\begin{equation}
  N_{\mathrm{dec}} \;\approx\;
   \log_{10}\!\left(\frac{L_{0}}{\eta}\right)
   \;\approx\; \log_{10}\!\left(\frac{0.75\,\mathrm{mm}}{7\mum}\right)
   \;\approx\; 2.0\,\mathrm{decadi}.
\end{equation}
Adottando i valori \emph{rivisti} secondo la critica \#1 (Pope
$\eta\approx 39\mum$):
\begin{equation}
  N_{\mathrm{dec},\mathrm{Pope}} \;\approx\;
   \log_{10}\!\left(\frac{0.75\,\mathrm{mm}}{39\mum}\right)
   \;\approx\; 1.3\,\mathrm{decadi}.
\end{equation}
Più realisticamente, una stima Tennekes--Lumley~\cite{Tennekes1972}
del range inerziale \emph{effettivamente sviluppato} per
$\Rey_{\lambda}=O(1)$ è di $\lesssim$ una mezza decade.

\subsection{Riformulazione operativa del criterio}

Il \texttt{main.pdf} ha \emph{già} qualificato il criterio Fase~4
come ``qualitativo'', non quantitativo. La forma operativa proposta
nel documento richiede:

\begin{enumerate}
  \item presenza di una zona di pendenza intermedia tra il range
    energetico (slope $\sim -1$) e il range dissipativo
    (slope $\sim -7$);
  \item slope locale $\in [-2.0, -1.3]$ in almeno mezza decade di
    frequenza tra $f \approx U/L_{0}$ e $f \approx 1/\tau_{\eta}$;
  \item assenza di \emph{pile-up} spettrale ad alte frequenze.
\end{enumerate}

Questi tre indicatori formano un test di plausibilità della
risoluzione DDES, non una richiesta di range inerziale esteso.

\subsection{Confronto con DNS di riferimento}

La DNS di Dairay et al.~\cite{Dairay2015} su getto impingente a
$\Rey=10\,000$ riporta circa 1.5 decadi di range inerziale. A
$\Rey=1\,000$ (caso Huawei), una scalatura
$N_{\mathrm{dec}}\propto \log\Rey^{3/4}$ predice
$\sim 0.6$ decadi. Questo è consistente con la stima
Tennekes--Lumley del paragrafo precedente.

\begin{verdettofalsa}
La critica ignora la nota esplicita della Sezione~8.3 del
\texttt{main.pdf} che declassa la verifica $-5/3$ a criterio
\emph{qualitativo}, con tolleranza esplicita per range inerziale
inferiore a una decade. Il \texttt{main.pdf} è anzi tra i pochi
documenti progettuali in letteratura ad anticipare e mitigare il
problema in modo trasparente: il criterio Fase~4 non è ``inarrivabile''
ma è stato deliberatamente formulato in forma morbida proprio per
tenere conto del basso $\Rey_{\lambda}$.
\end{verdettofalsa}

% ============================================================
\section{Critica \#5 — Inlet Plug Flow Senza Turbolenza Sintetica}
\label{sec:c5}
% ============================================================

\begin{citazione}
\textit{L'inlet è troppo pulito. L'uso di un profilo plug flow
uniforme senza turbolenza sintetica (Opzione A) per un tubo così
corto ritarderà artificialmente lo sviluppo della turbolenza nello
shear layer, alterando severamente la previsione del rumore
quadrupolare.}
\end{citazione}

\subsection{Cosa dice il \texttt{main.pdf}}

La Sezione~10.4 del \texttt{main.pdf} (\emph{Decisioni Aperte ---
Turbolenza di Ingresso}) riconosce \emph{esplicitamente} il problema:

\begin{quote}\small\itshape
``Il profilo di velocità uniforme all'inlet non genera fluttuazioni
di turbolenza sintetiche. In un getto reale, la turbolenza residua
nell'ugello eccita più rapidamente l'instabilità di Kelvin--Helmholtz
nello shear layer e accelera la transizione.''
\end{quote}

Il documento offre due opzioni:
\begin{description}
  \item[Opzione A:] inlet plug flow deterministico, riproducibile.
  \item[Opzione B:] inlet stocastico con metodo di Klein
    (\ttcfg{INLET\_TURBULENCE= YES}, intensità e scala
    specificate).
\end{description}

\subsection{Sviluppo della shear layer di Kelvin--Helmholtz}

Il numero di Strouhal della modalità più amplificata di
Kelvin--Helmholtz su un getto piano a profilo top-hat è
$\Str_{\mathrm{KH}}\approx 0.017$~\cite{Michalke1965}, con tasso di
crescita $\sigma_{\mathrm{KH}}/U_{\mathrm{jet}}\approx 0.18/\Djet$.
Per il caso Huawei, il tempo di e-folding di KH è:
\begin{equation}
  \tau_{\mathrm{KH}}^{(e)} \;=\; \frac{\Djet}{0.18\,U}
   \;=\; \frac{1.5\times 10^{-3}}{0.18\times 10}
   \;\approx\; 833\mus.
\end{equation}

Il tempo di transito attraverso il tubo inlet ($H_{\mathrm{tube}}=
1\,\mathrm{mm}$) è soltanto:
\begin{equation}
  t_{\mathrm{tube}} \;=\; \frac{H_{\mathrm{tube}}}{U}
   \;=\; \frac{10^{-3}}{10}
   \;=\; 100\mus.
\end{equation}

Il rapporto $t_{\mathrm{tube}}/\tau_{\mathrm{KH}}^{(e)}\approx 0.12$
indica che, anche con turbolenza sintetica iniettata, la
perturbazione \emph{non avrebbe il tempo di amplificarsi
significativamente} prima dell'impatto sul chip. Il tubo è
elasticamente troppo corto.

\subsection{Implicazioni quantitative sul rumore}

Il rumore aerodinamico di un getto si compone di due contributi
(Lighthill 1952, Goldstein 1976):
\begin{itemize}
  \item \textbf{Sorgenti dipolari} (parete): dominanti a basso
    Mach. SPL $\propto M^{6}$.
  \item \textbf{Sorgenti quadrupolari} (volume): SPL $\propto M^{8}$.
\end{itemize}

Per il caso Huawei con $\Ma=0.029$, il rapporto quadrupolari/dipolari
è:
\begin{equation}
  \frac{\mathrm{SPL}_{Q}}{\mathrm{SPL}_{D}} \;\propto\; \Ma^{2}
   \;=\; (0.029)^{2} \;\approx\; 8.4\times 10^{-4}.
\end{equation}

Le sorgenti quadrupolari sono dunque circa $30\,\mathrm{dB}$ sotto
le dipolari, e \emph{in questo regime di basso Mach il rumore
totale è dominato quasi completamente dalle sorgenti dipolari di
parete}. La preoccupazione del peer-reviewer riguardo al ``rumore
quadrupolare'' è quindi quantitativamente di importanza secondaria
per il caso specifico.

\subsection{Effetto sulla validazione PIV}

Il confronto principale con i dati sperimentali Huawei è basato su
profili di velocità media (PIV) e pressione superficiale (microfono
sul chip). Per questi confronti, la condizione di inlet
\emph{deterministica} è preferibile perché:

\begin{itemize}
  \item Le condizioni sperimentali Huawei usano un compressore
    laminare con calmiere a nido d'ape: turbolenza di inlet
    misurata $<1\%$ ($u'_{\mathrm{inlet}}/U \approx 0.005$).
  \item Aggiungere turbolenza sintetica con $u'/U=0.05$--$0.10$
    (valori tipici dei generatori) introdurrebbe una
    \emph{sovrastima} della turbolenza di inlet.
  \item L'incertezza sui parametri della turbolenza sintetica
    (intensità, scala integrale, tensore di Reynolds) è una sorgente
    di errore di validazione difficile da quantificare.
\end{itemize}

\subsection{Quando l'Opzione B diventa necessaria}

L'Opzione B (turbolenza sintetica) è raccomandata in due scenari:
\begin{enumerate}
  \item Studio di sensitivity per quantificare l'errore della
    Opzione A: stessa simulazione ripetuta con TI=$1\%$ e
    $L_{t}=0.1\,\Djet$, confronto dei profili di velocità a
    $z/H=0.5$.
  \item Configurazione futura con tubo inlet più lungo
    ($H_{\mathrm{tube}}\geq 5\,\Djet$), dove la transizione naturale
    del flusso ha tempo di svilupparsi.
\end{enumerate}

\subsection{Azione correttiva proposta}

\begin{infobox}[Studio di sensitivity inlet]
A valle della prima campagna WRLES con Opzione A, eseguire
\textbf{una singola simulazione di sensitivity} (durata 4 settimane,
costo $\sim$\texteuro 400) con Opzione B
(\ttcfg{INLET\_TURBULENCE= YES}, $\mathrm{TI}=1\%$,
$L_{t}=0.1\Djet$) e confrontare:
\begin{itemize}
  \item profilo di velocità media $\bar u_{z}(r)$ a $z/H=0.5$,
  \item RMS $u'_{z}(r)$ nello shear layer,
  \item SPL al microfono sul chip.
\end{itemize}
Se le differenze sono $<5\%$ sui profili medi e $<2\,\mathrm{dB}$
sull'SPL, la scelta dell'Opzione A è validata. In caso contrario,
adottare l'Opzione B per il caso principale.
\end{infobox}

\begin{verdettoparz}
La critica è \emph{fisicamente corretta} sull'osservazione che la
turbolenza di inlet accelera la transizione di Kelvin--Helmholtz.
Tuttavia: (a) il tubo inlet del caso Huawei è troppo corto perché
una perturbazione di inlet abbia il tempo di amplificarsi
significativamente ($t_{\mathrm{tube}}/\tau_{\mathrm{KH}}^{(e)}\approx
0.12$); (b) le sorgenti quadrupolari sono $\sim 30\,\mathrm{dB}$
sotto le dipolari nel regime $\Ma=0.029$, quindi il rumore non è
dominato da sorgenti volumetriche; (c) la condizione sperimentale
Huawei usa flusso di inlet quasi-laminare ($\mathrm{TI}<1\%$),
rendendo l'Opzione A più rappresentativa della realtà. Lo studio di
sensitivity con Opzione B (suggerito nello schedulato di rischio
Sezione~10.4 del \texttt{main.pdf}) resta consigliato come
verifica.
\end{verdettoparz}

% ============================================================
\section{Conclusioni}
\label{sec:concl}
% ============================================================

L'analisi quantitativa delle cinque critiche conduce ai seguenti
punti di sintesi:

\begin{enumerate}
  \item Le critiche \#1, \#2, \#5 sono \emph{parzialmente fondate}:
    sollevano problemi reali ma di entità contenuta o
    consapevolmente accettati nel \texttt{main.pdf}. Per ognuna
    è proposta un'azione correttiva implementabile in fase
    iterativa, senza modificare la struttura della pipeline.

  \item Le critiche \#3 e \#4 sono \emph{infondate}: la prima
    confonde il parametro \texttt{JST\_SENSOR\_COEFF} con uno
    schema di tipo JST, ignorando che il setup
    $\mathrm{FDS}+\mathrm{MUSCL}+(\kappa_{2}=0,\kappa_{4}=0.02)$ è
    la best-practice SU2 v8 per LES incompressibili; la seconda
    ignora il declassamento esplicito a criterio \emph{qualitativo}
    già presente nella Sezione~8.3 del \texttt{main.pdf}.

  \item Nessuna delle critiche, sole o combinate, mette in
    discussione la \emph{fondatezza scientifica} della pipeline
    multi-fedeltà RANS$\to$URANS$\to$DDES$\to$WMLES$\to$WRLES$\to$FW-H.
    Le azioni correttive proposte (Sezioni~\ref{sec:c1},~\ref{sec:c2},
    \ref{sec:c5}) sono \emph{raffinamenti incrementali},
    implementabili durante l'esecuzione della commessa senza
    impatto sulla roadmap a 4 mesi.

  \item Si propone di \emph{integrare} formalmente le tre azioni
    correttive nella prossima revisione del \texttt{main.pdf}
    (versione 1.1):
    \begin{itemize}
      \item Aggiornare la Sezione~6.2 con il commento
        sull'\emph{aggiornamento iterativo} di $\varepsilon$ basato
        sulla WMLES L4.
      \item Aggiungere alla Sezione~7.3 la formula di Sibulkin per
        $C_{f,\mathrm{stag}}$ e l'indicazione di raffinamento
        locale nella zona di stagnazione.
      \item Promuovere il \emph{sensitivity check inlet} (Opzione B)
        da rischio mitigabile (Sezione~10.4) a milestone
        intermedia (M2.5) della roadmap.
    \end{itemize}
\end{enumerate}

Il valore del peer-review esterno risiede precisamente nel forzare
la verifica esplicita di assunzioni che, se accettate
implicitamente, possono propagarsi attraverso l'intera pipeline.
Le tre obiezioni parzialmente fondate hanno permesso di identificare
altrettanti raffinamenti operativi a costo modesto, mentre le due
critiche infondate hanno offerto l'occasione di documentare con
maggiore chiarezza scelte numeriche già consapevolmente compiute.

\bigskip
\centerline{\itshape\color{midblue}--- Fine del documento ---}

% ============================================================
% RIFERIMENTI
% ============================================================
\newpage
\begin{thebibliography}{99}
\small

\bibitem{Pope2000}
  Pope, S.B. (2000).
  \textit{Turbulent Flows}. Cambridge University Press.

\bibitem{Sibulkin1952}
  Sibulkin, M. (1952).
  Heat transfer near the forward stagnation point of a body of revolution.
  \textit{Journal of Aeronautical Sciences}, 19(8), 570--571.

\bibitem{Hadziabdic2008}
  Hadziabdic, M. \& Hanjalic, K. (2008).
  Vortical structures and heat transfer in a round impinging jet.
  \textit{Journal of Fluid Mechanics}, 596, 221--260.

\bibitem{Dairay2015}
  Dairay, T., Fortune, V., Lamballais, E. \& Brizzi, L.-E. (2015).
  Direct numerical simulation of a turbulent jet impinging on a heated wall.
  \textit{Journal of Fluid Mechanics}, 764, 362--394.

\bibitem{Economon2016}
  Economon, T.D., Palacios, F., Copeland, S.R., Lukaczyk, T.W. \& Alonso, J.J. (2016).
  SU2: An Open-Source Suite for Multiphysics Design and Analysis.
  \textit{AIAA Journal}, 54(3), 828--846.

\bibitem{Aupoix2017}
  Aupoix, B. \& Spalart, P.R. (2017).
  Extensions of the Spalart-Allmaras turbulence model to account for
  wall roughness.
  \textit{International Journal of Heat and Fluid Flow}, 24(4), 454--462.

\bibitem{Tennekes1972}
  Tennekes, H. \& Lumley, J.L. (1972).
  \textit{A First Course in Turbulence}. MIT Press.

\bibitem{Michalke1965}
  Michalke, A. (1965).
  On spatially growing disturbances in an inviscid shear layer.
  \textit{Journal of Fluid Mechanics}, 23(3), 521--544.

\bibitem{Spalart2006}
  Spalart, P.R., Deck, S., Shur, M.L., Squires, K.D., Strelets, M.Kh. \& Travin, A. (2006).
  A new version of detached-eddy simulation, resistant to ambiguous grid densities.
  \textit{Theoretical and Computational Fluid Dynamics}, 20(3), 181--195.

\bibitem{Celik2008}
  Celik, I.B., Ghia, U., Roache, P.J., Freitas, C.J., Coleman, H., Raad, P.E. (2008).
  Procedure for Estimation and Reporting of Uncertainty Due to Discretization in CFD Applications.
  \textit{Journal of Fluids Engineering}, 130(7), 078001.

\end{thebibliography}

\end{document}
